% 22
\begin{frame}{Касательная к выпуклой функции}
\centering
\includegraphics[width=0.79\linewidth]{tangent_lower_bound.pdf}

\vspace{0.2em}
\small
В отличие от обычной локальной аппроксимации, у выпуклой функции касательная оказывается \textbf{глобальной нижней оценкой}.
\end{frame}

% 23
\begin{frame}{Критерий первого порядка}
\small
\begin{block}{Теорема}
Пусть $f$ дифференцируема на выпуклом множестве $X$. Тогда
\[
\boxed{
f\text{ выпукла}
\iff
f(y)\ge f(x)+\grad f(x)^\top(y-x)
}
\]
для любых $x,y\in X$.
\end{block}

Правая часть --- касательная гиперплоскость к графику в точке $x$.
\end{frame}

% 24
\begin{frame}{Связь с локальной моделью из лекции 2}
\small
Раньше мы писали
\[
 f(x+h)\approx f(x)+\grad f(x)^\top h.
\]
Это была \textbf{локальная} линейная модель.

Для выпуклой функции при $h=y-x$ получаем уже точное глобальное утверждение:
\[
\boxed{
 f(y)\ge f(x)+\grad f(x)^\top(y-x).
}
\]

\begin{alertblock}{Главная мысль}
Линейная модель выпуклой функции никогда не переоценивает саму функцию.
\end{alertblock}
\end{frame}

% 25
\begin{frame}{Доказательство критерия: направление $\Rightarrow$}
\footnotesize
Предположим, что $f$ выпукла. Зафиксируем $x,y$ и определим функцию одной переменной
\[
 \varphi(t)=f(x+t(y-x)),\qquad t\in[0,1].
\]
По выпуклости
\[
 \varphi(t)\le(1-t)\varphi(0)+t\varphi(1).
\]
Переносим $\varphi(0)$ и делим на $t>0$:
\[
 \frac{\varphi(t)-\varphi(0)}{t}
 \le
 \varphi(1)-\varphi(0).
\]
\end{frame}

% 26
\begin{frame}{Доказательство критерия: предел при $t\to0+$}
\footnotesize
Переходим к пределу:
\[
 \varphi'(0)
 \le
 f(y)-f(x).
\]
Но по правилу цепочки
\[
 \varphi'(0)
 =\grad f(x)^\top(y-x).
\]
Следовательно,
\[
\boxed{
 f(y)\ge f(x)+\grad f(x)^\top(y-x).
}
\]
\end{frame}

% 27
\begin{frame}{Обратное направление критерия}
\footnotesize
Теперь предположим, что для любых $u,v$
\[
 f(v)\ge f(u)+\grad f(u)^\top(v-u).
\]
Положим
\[
 z=(1-t)x+ty.
\]
Применим условие в точке $z$ отдельно к $x$ и $y$:
\[
 f(x)\ge f(z)+\grad f(z)^\top(x-z),
\]
\[
 f(y)\ge f(z)+\grad f(z)^\top(y-z).
\]
\end{frame}

% 28
\begin{frame}{Завершение обратного направления}
\footnotesize
Умножим первое неравенство на $1-t$, второе на $t$ и сложим:
\[
\begin{aligned}
(1-t)f(x)+tf(y)
&\ge f(z)\\
&\quad+\grad f(z)^\top\big((1-t)(x-z)+t(y-z)\big).
\end{aligned}
\]
Но
\[
(1-t)(x-z)+t(y-z)=0.
\]
Поэтому
\[
 f((1-t)x+ty)\le(1-t)f(x)+tf(y).
\]
Это и есть выпуклость.
\end{frame}

% 29
\begin{frame}{Что стоит запомнить из критерия первого порядка}
\small
Не обязательно начинать с формулы. Геометрически утверждение звучит так:
\[
\boxed{\text{касательная к выпуклой функции — глобальная нижняя оценка}.}
\]

Это позволяет превращать локальную информацию $\grad f(x)$ в утверждение о функции во всём пространстве.

\begin{alertblock}{Для коллоквиума}
Важно уметь и сформулировать критерий, и восстановить доказательство через функцию $\varphi(t)=f(x+t(y-x))$.
\end{alertblock}
\end{frame}

% 30
\begin{frame}{Самопроверка: квадратичная функция}
\small
Пусть
\[
 f(x)=\frac12x^\top Ax,
 \qquad A=A^\top.
\]
Критерий первого порядка требует
\[
 \frac12y^\top Ay
 \ge
 \frac12x^\top Ax+x^\top A(y-x).
\]
После упрощения:
\[
\boxed{
\frac12(y-x)^\top A(y-x)\ge0.
}
\]

Какое условие на $A$ обеспечивает это для любых $x,y$?
\end{frame}

% 31
\begin{frame}{Локальный минимум выпуклой функции — глобальный}
\small
\begin{block}{Теорема}
Пусть $f$ выпукла на выпуклом множестве $X$. Тогда любой локальный минимум $x^\star$ функции $f$ является глобальным.
\end{block}

\centering
\includegraphics[width=0.78\linewidth]{local_global_minima.pdf}
\end{frame}

% 32
\begin{frame}{Доказательство: предположим обратное}
\small
Пусть $x^\star$ --- локальный минимум, но существует $y\in X$ такое, что
\[
 f(y)<f(x^\star).
\]
Рассмотрим точки
\[
 x_t=(1-t)x^\star+ty,
 \qquad t\in(0,1).
\]
При малом $t>0$ точка $x_t$ находится сколь угодно близко к $x^\star$.
\end{frame}

% 33
\begin{frame}{Доказательство: противоречие с локальностью}
\small
По выпуклости
\[
\begin{aligned}
f(x_t)
&\le(1-t)f(x^\star)+tf(y)\\
&< (1-t)f(x^\star)+tf(x^\star)\\
&=f(x^\star).
\end{aligned}
\]

Но $x_t$ можно выбрать сколь угодно близко к $x^\star$. Это противоречит тому, что $x^\star$ --- локальный минимум.
\end{frame}

% 34
\begin{frame}{Условие оптимальности для дифференцируемой выпуклой функции}
\small
\begin{block}{Теорема}
Если $f$ дифференцируема и выпукла на $\R^d$, то
\[
\boxed{
 \grad f(x^\star)=0
 \iff
 x^\star\text{ — глобальный минимум }f.
}
\]
\end{block}

Левая часть --- локальное условие. Выпуклость превращает его в глобальный результат.
\end{frame}

% 35
\begin{frame}{Почему $\grad f(x^\star)=0$ достаточно}
\small
Из критерия первого порядка для любого $y$:
\[
 f(y)
 \ge
 f(x^\star)+\grad f(x^\star)^\top(y-x^\star).
\]
Если
\[
 \grad f(x^\star)=0,
\]
то
\[
 \boxed{f(y)\ge f(x^\star)\quad\forall y.}
\]
Следовательно, $x^\star$ --- глобальный минимум.
\end{frame}

% 36
\begin{frame}{Почему это фундаментально}
\small
Для произвольной гладкой функции условие
\[
 \grad f(x)=0
\]
может означать:
\[
\text{минимум},\qquad
\text{максимум},\qquad
\text{седловую точку}.
\]

Для дифференцируемой выпуклой функции:
\[
\boxed{\grad f(x)=0\Longrightarrow\text{глобальный минимум}.}
\]

Это одна из главных причин, почему выпуклые задачи настолько удобны.
\end{frame}
